\chapter{Preface}
\label{ch:preface}

This book began as a study journal. I run Spark SQL on an Apache Iceberg
lakehouse every working day --- Spark 3.5.2, Iceberg v2, Kyuubi in front, Hive
Metastore as the catalog, Google Cloud Storage underneath, all on GKE --- and I
kept noticing the same thing about how I, and the people around me, fixed slow
queries. We changed a setting, reran the query, kept the change if it helped, and
moved on without knowing why it had. That works while a query takes a minute.
It stops working when each trial takes an hour and costs real money, and it stops
working entirely when the setting that helped last quarter makes this quarter's
query slower.

The notes that became this book were an attempt to replace that trial loop with a
model. The model is one sentence, and the rest of the book is its consequences:

\begin{quote}
\itshape
Every Spark performance decision is a choice about where and when data moves ---
between operators, across the network, and between metadata and storage --- and
you can predict and change those movements once you can see them.
\end{quote}

A scan reads files; a shuffle moves rows between executors so that rows with the
same key meet; a join decides which side to move; adaptive execution changes its
mind once it has seen real sizes; Iceberg's planner reads a few kilobytes of
metadata so it does not have to read gigabytes of data. These look like separate
subjects. They are one subject, and once you see them as one, the question you
ask of a slow query changes from ``which config should I try?'' to ``what is
moving, how much, and does it have to?''

\section*{What makes this book different}

There are good books on Spark and good material on Iceberg. This one is organized
around two commitments.

\paragraph{Mechanism before settings.} This is not a list of
\texttt{spark.sql.*} keys with recommended values. A configuration appears only
after the mechanism it controls has been explained, and always with what it
costs. A reader who finishes should be able to derive most of the ``recommended''
values from first principles, and to recognize the case where the recommendation
is wrong. When a rule of thumb is useful --- and some are --- the book shows where
it comes from.

\paragraph{The seam between Catalyst and Iceberg.} Most of the cost of a lakehouse
query is decided at the boundary between the engine and the table format: what the
planner can push into the scan, what the manifests let it skip, how the file
layout written last night shapes the tasks run this morning. Treating Spark and
Iceberg as two separate subjects hides that boundary. This book is built around
it: Part~II ends with pruning and bucketing on the Spark side, and Part~III picks
the same threads up on the Iceberg side.

\section*{A running specimen, and numbers you can reproduce}

The book follows one query, introduced in \Cref{ch:shape}: a week of orders joined
to a product dimension and aggregated to an on-time delivery rate. On the
production-sized cluster it takes about four minutes. We do not try to make it
faster in the first chapter; we learn to \emph{read} it --- the plan, the stages,
the shuffle, the one task that runs longer than the rest --- and then every
chapter changes one thing about it and shows the consequence.

Wherever I could, I measured rather than asserted. The experiments run in a local
Docker container with Spark 3.5.2 and Iceberg 1.6.1 against synthetic data
generated to the same shape as the production tables (\Cref{sec:refenv-lab}), and
the scripts live alongside the manuscript. Several times the measurement
contradicted what I had written from memory: a default that I believed applied
did not fire, a threshold compared a different quantity than I had assumed, a
mechanism needed a second setting before it did anything. Those corrections are
in the text, where they belong, because the gap between what a setting is
documented to do and what it does on your data is precisely where tuning goes
wrong.

Some claims I have not tested. These are marked ``claim --- not tested'' (or
``verify on your cluster'') rather than smoothed over. Treat them as hypotheses.

\section*{How the book is organized}

\textbf{Part~I, How Spark Executes a Query,} builds the machinery that moves data
and teaches you to see it: the driver, executors, jobs, stages and tasks
(\Cref{ch:shape}); the Catalyst pipeline from SQL text to physical plan
(\Cref{ch:catalyst}); and the executor's memory and CPU budget
(\Cref{ch:execmem}).

\textbf{Part~II, Making Queries Fast,} treats each chapter as one lever on data
movement: the shuffle (\Cref{ch:shuffle}), joins (\Cref{ch:joins}), adaptive
query execution (\Cref{ch:aqe}), pruning, partitioning, and bucketing
(\Cref{ch:pruning}), and caching and reuse (\Cref{ch:caching}).

\textbf{Part~III, The Iceberg Lakehouse,} extends the model below the file:
metadata is data that moves too. It covers the table format itself
(\Cref{ch:iceberg-spec}), catalogs and the commit protocol
(\Cref{ch:catalogs}), partitioning and layout (\Cref{ch:iceberg-layout}), reading
efficiently (\Cref{ch:iceberg-read}), and writing and maintaining tables
(\Cref{ch:iceberg-write}).

\textbf{Part~IV, Operating Spark,} applies the model under real constraints:
sizing and configuration (\Cref{ch:tuning}), debugging and profiling
(\Cref{ch:debugging}), and streaming into the lakehouse (\Cref{ch:streaming}).

\textbf{Part~V, Internals,} goes one level deeper on the same machinery: Tungsten
and code generation (\Cref{ch:tungsten}), the shuffle framework and scheduler
(\Cref{ch:scheduler}), and DataSource~V2 (\Cref{ch:dsv2}). It is optional, and the
earlier parts do not depend on it.

Five appendices follow: the reference environment and running example
(\Cref{app:refenv}), a configuration reference (\Cref{app:config}), a guide to
reading query plans (\Cref{app:plans}), a glossary (\Cref{app:glossary}), and
suggestions for further reading (\Cref{app:reading}).

\section*{A note on scope}

The book is written against one stack, and it says so every time a claim depends
on it. Spark and Iceberg both change quickly, and behavior that is true on 3.5.2
may not be true on 4.0. I would rather give you a model you can re-test against
your version than a table of values that quietly goes stale.

\section*{Acknowledgments}

This book stands on the Apache Spark and Apache Iceberg projects --- their source
code, their specifications, and the documentation and talks of the people who
maintain them. Every number measured here was possible because those projects are
open and can be run on a laptop. Errors that remain are mine, and corrections are
welcome through the repository.

\bigskip
\begin{flushright}
\textit{Alan Phan}\\
2026
\end{flushright}
